% answers.tex —— 2014 年·安徽卷·理科 · 参考答案与详细解析片段
% 由 tools/generate_projects.py 依来源总账自 2014/2014_数学真题_安徽理_试题.md + 2014/2014_数学真题_安徽理_解析正文.md + 2014/2014_数学真题_安徽理_答案要点.md 生成，勿手改（改源 MD 后重跑）


\begin{answer}
\gksolhead{1}{}

\textbf{题目}：

设 $\mathrm i$ 是虚数单位，$\bar z$ 表示复数 $z$ 的共轭复数。若 $z=1+\mathrm i$，则 $\dfrac{z}{\mathrm i}+\mathrm i\cdot\bar z=$\mbox{（　　）}

（A）$-2$；（B）$-2\mathrm i$；（C）$2$；（D）$2\mathrm i$。

\Needspace{4\baselineskip}
\textbf{解析}

\Needspace{5\baselineskip}
由 $z=1+\mathrm i$ 得 $\bar z=1-\mathrm i$，且 $\dfrac1{\mathrm i}=-\mathrm i$。分别计算两项：
\nopagebreak[4]
\[\frac{z}{\mathrm i}=(1+\mathrm i)(-\mathrm i)=1-\mathrm i,\qquad
\mathrm i\cdot\bar z=\mathrm i(1-\mathrm i)=1+\mathrm i .\]
\Needspace{5\baselineskip}
相加得
\nopagebreak[4]
\[\frac{z}{\mathrm i}+\mathrm i\cdot\bar z=(1-\mathrm i)+(1+\mathrm i)=2 .\]

选 C。

\gkanswer{答案：C。$\dfrac{z}{\mathrm i}=1-\mathrm i$，$\mathrm i\bar z=1+\mathrm i$，相加得 $2$。}

\end{answer}
\begin{answer}
\gksolhead{2}{}

\textbf{题目}：

“$x<0$”是“$\ln(x+1)<0$”的\mbox{（　　）}

（A）充分不必要条件；（B）必要不充分条件；（C）充分必要条件；（D）既不充分也不必要条件。

\Needspace{4\baselineskip}
\textbf{解析}

\Needspace{5\baselineskip}
先解对数不等式。由定义域要求 $x+1>0$，即 $x>-1$。在此前提下
\nopagebreak[4]
\[\ln(x+1)<0\iff 0<x+1<1\iff -1<x<0 .\]
因此“$\ln(x+1)<0$”推出“$x<0$”，反过来不成立：例如 $x=-2$ 满足 $x<0$，但 $x+1=-1$ 使对数无定义。所以“$x<0$”是“$\ln(x+1)<0$”的必要不充分条件。

选 B。

\gkanswer{答案：B。$\ln(x+1)<0\iff-1<x<0$，故“$x<0$”必要但不充分。}

\end{answer}
\begin{answer}
\gksolhead{3}{}

\textbf{题目}：

如图所示，程序框图（算法流程图）的输出结果是\mbox{（　　）}

（A）$34$；（B）$55$；（C）$78$；（D）$89$。
\nopagebreak[4]
\begin{center}% ============================================================
% f2014_anhui_science_q03_flowchart.tex
% 2014 安徽卷（理科）第 3 题：程序框图
% 来源 PDF：普通高考/2014/2014安徽理.pdf
% 原页：第 1 页（A3 横排）左栏，印面 857，第 3 题题干下方
% 结构：开始 → 赋初值 x=1、y=1 → z=x+y → 判断 z≤50；
%       是：x=y → y=z → 返回 z=x+y 之前；否：输出 z → 结束。
% ============================================================
\begin{tikzpicture}[>=Stealth, line width=0.7pt, font=\small,
  x=1cm, y=1cm, line cap=round, line join=round]

  % 开始
  \draw[rounded corners=6pt] (-0.95,0.42) rectangle (0.95,-0.42);
  \node at (0,0) {开始};
  \draw[->] (0,-0.42) -- (0,-0.85);

  % 赋初值
  \draw (-1.55,-0.85) rectangle (1.55,-1.70);
  \node at (0,-1.28) {$x=1$，$y=1$};
  \draw[->] (0,-1.70) -- (0,-2.10);

  % z=x+y
  \draw (-1.30,-2.10) rectangle (1.30,-2.95);
  \node at (0,-2.52) {$z=x+y$};
  \draw[->] (0,-2.95) -- (0,-3.40);

  % 判断框
  \draw (0,-3.40) -- (1.80,-4.30) -- (0,-5.20) -- (-1.80,-4.30) -- cycle;
  \node at (0,-4.30) {$z\leqslant 50$};

  % 是：向下进入循环体
  \draw[->] (0,-5.20) -- (0,-5.60);
  \node at (0.58,-5.42) {是};
  \draw (-1.20,-5.60) rectangle (1.20,-6.45);
  \node at (0,-6.02) {$x=y$};

  \draw[->] (0,-6.45) -- (0,-6.85);
  \draw (-1.20,-6.85) rectangle (1.20,-7.70);
  \node at (0,-7.27) {$y=z$};

  % 回边：向下、向左、向上，再水平汇入 z=x+y 上方
  \draw (-1.20,-7.27) -- (-2.95,-7.27) -- (-2.95,-1.90) -- (0,-1.90);
  \draw[->] (0,-1.90) -- (0,-2.10);

  % 否：向右输出
  \draw[->] (1.80,-4.30) -- (2.98,-4.30);
  \node at (2.35,-3.98) {否};
  \draw (3.15,-3.88) -- (5.35,-3.88) -- (5.05,-4.72) -- (2.85,-4.72) -- cycle;
  \node at (4.10,-4.30) {输出 $z$};
  \draw[->] (4.10,-4.72) -- (4.10,-5.10);

  % 结束
  \draw[rounded corners=6pt] (3.15,-5.10) rectangle (5.05,-5.94);
  \node at (4.10,-5.52) {结束};

\end{tikzpicture}
\end{center}
\nopagebreak[4]
\Needspace{4\baselineskip}
\textbf{解析}

初值 $x=1$，$y=1$。每轮先算 $z=x+y$，若 $z\leqslant50$ 则令 $x$ 取旧的 $y$、$y$ 取 $z$，再返回计算下一轮的 $z$；否则输出 $z$。

\begin{center}
\begin{tabular}{llllll}
\toprule
轮次 & 计算前 $x$ & 计算前 $y$ & $z=x+y$ & $z\leqslant50$ & 新 $(x,y)$ \\
\midrule
1 & $1$ & $1$ & $2$ & 是 & $(1,2)$ \\
2 & $1$ & $2$ & $3$ & 是 & $(2,3)$ \\
3 & $2$ & $3$ & $5$ & 是 & $(3,5)$ \\
4 & $3$ & $5$ & $8$ & 是 & $(5,8)$ \\
5 & $5$ & $8$ & $13$ & 是 & $(8,13)$ \\
6 & $8$ & $13$ & $21$ & 是 & $(13,21)$ \\
7 & $13$ & $21$ & $34$ & 是 & $(21,34)$ \\
8 & $21$ & $34$ & $55$ & 否 & 输出 $55$ \\
\bottomrule
\end{tabular}
\end{center}

输出的 $z$ 依次为 $2,3,5,8,13,21,34,55$，第一个大于 $50$ 的是 $55$，故输出 $55$。

选 B。

\gkanswer{答案：B。$z$ 依次取 $2,3,5,8,13,21,34,55$，第一个大于 $50$ 的是 $55$。}

\end{answer}
\begin{answer}
\gksolhead{4}{}

\textbf{题目}：

以平面直角坐标系的原点为极点，$x$ 轴的正半轴为极轴，建立极坐标系，两种坐标系中取相同的长度单位。已知直线 $l$ 的参数方程是 $\begin{cases}x=t+1,\\y=t-3,\end{cases}$（$t$ 为参数），圆 $C$ 的极坐标方程是 $\rho=4\cos\theta$，则直线 $l$ 被圆 $C$ 截得的弦长为\mbox{（　　）}

（A）$\sqrt{14}$；（B）$2\sqrt{14}$；（C）$\sqrt2$；（D）$2\sqrt2$。

\Needspace{4\baselineskip}
\textbf{解析}

\Needspace{5\baselineskip}
由参数方程消去 $t$：$t=x-1$，代入 $y=t-3$ 得
\nopagebreak[4]
\[y=x-4,\qquad\text{\text{即}}\quad x-y-4=0 .\]
\Needspace{5\baselineskip}
由 $\rho=4\cos\theta$ 得 $\rho^2=4\rho\cos\theta$，即
\nopagebreak[4]
\[x^2+y^2=4x\iff (x-2)^2+y^2=4,\]
圆心为 $(2,0)$，半径 $r=2$。

\Needspace{5\baselineskip}
圆心到直线 $x-y-4=0$ 的距离
\nopagebreak[4]
\[d=\frac{|2-0-4|}{\sqrt{1^2+(-1)^2}}=\frac2{\sqrt2}=\sqrt2 .\]
\Needspace{5\baselineskip}
故弦长为
\nopagebreak[4]
\[2\sqrt{r^2-d^2}=2\sqrt{4-2}=2\sqrt2 .\]

选 D。

\gkanswer{答案：D。直线为 $x-y-4=0$，圆心 $(2,0)$ 到它的距离为 $\sqrt2$，半径 $2$，弦长 $2\sqrt2$。}

\end{answer}
\begin{answer}
\gksolhead{5}{}

\textbf{题目}：

$x$，$y$ 满足约束条件 $\begin{cases}x+y-2\leqslant0,\\x-2y-2\leqslant0,\\2x-y+2\geqslant0,\end{cases}$ 若 $z=y-ax$ 取得最大值的最优解不唯一，则实数 $a$ 的值为\mbox{（　　）}

（A）$\dfrac12$ 或 $-1$；（B）$2$ 或 $\dfrac12$；（C）$2$ 或 $1$；（D）$2$ 或 $-1$。

\Needspace{4\baselineskip}
\textbf{解析}

\Needspace{5\baselineskip}
三条边界两两相交，得可行域的顶点：
\nopagebreak[4]
\[x+y=2\ \text{\text{与}}\ x-2y=2\ \Rightarrow\ (2,0),\]
\nopagebreak[4]
\[x+y=2\ \text{\text{与}}\ 2x-y=-2\ \Rightarrow\ (0,2),\]
\nopagebreak[4]
\[x-2y=2\ \text{\text{与}}\ 2x-y=-2\ \Rightarrow\ (-2,-2).\]
目标函数 $z=y-ax$ 的等值线为 $y=ax+z$。最优解不唯一，必须有一条边界与等值线平行且整条边都取得最大值。

\Needspace{5\baselineskip}
当 $a=2$ 时，$z=y-2x$，三个顶点处的值依次为
\nopagebreak[4]
\[(2,0):-4,\qquad (0,2):2,\qquad (-2,-2):2 .\]
顶点 $(0,2)$ 与 $(-2,-2)$ 连成的边界上 $z$ 恒等于 $2$，且其余顶点处更小，故最大值在整条边上取得。

\Needspace{5\baselineskip}
当 $a=-1$ 时，$z=y+x$，三个顶点处的值依次为
\nopagebreak[4]
\[(2,0):2,\qquad (0,2):2,\qquad (-2,-2):-4 .\]
顶点 $(2,0)$ 与 $(0,2)$ 连成的边界上 $z$ 恒等于 $2$，最大值也在整条边上取得。

等值线与某条边界平行时，由三条边界的斜率可得 $a=2$，$a=-1$ 或 $a=\dfrac12$。当 $a=\dfrac12$ 时，边 $(2,0)$—$(-2,-2)$ 上 $z\equiv-1$，但顶点 $(0,2)$ 处 $z=2$，该边是最小值边而非最大值边，最优解仍唯一。因此 $a=2$ 或 $a=-1$。

选 D。

\gkanswer{答案：D。$a=2$ 时边 $(0,2)$—$(-2,-2)$ 上 $z\equiv2$ 为最大值；$a=-1$ 时边 $(2,0)$—$(0,2)$ 上 $z\equiv2$ 为最大值；$a=\dfrac12$ 时对应边为最小值边，不合。}

\end{answer}
\begin{answer}
\gksolhead{6}{}

\textbf{题目}：

设函数 $f(x)$（$x\in\mathbf R$）满足 $f(x+\pi)=f(x)+\sin x$。当 $0\leqslant x<\pi$ 时，$f(x)=0$，则 $f\left(\dfrac{23\pi}6\right)=$\mbox{（　　）}

（A）$\dfrac12$；（B）$\dfrac{\sqrt3}2$；（C）$0$；（D）$-\dfrac12$。

\Needspace{4\baselineskip}
\textbf{解析}

\Needspace{5\baselineskip}
由 $f(x+\pi)=f(x)+\sin x$ 逐次推进：
\nopagebreak[4]
\[f(x+2\pi)=f(x+\pi)+\sin(x+\pi)=f(x)+\sin x-\sin x=f(x),\]
\nopagebreak[4]
\[f(x+3\pi)=f(x+2\pi)+\sin(x+2\pi)=f(x)+\sin x .\]
\Needspace{5\baselineskip}
又
\nopagebreak[4]
\[\frac{23\pi}6=3\pi+\frac{5\pi}6,\]
\Needspace{5\baselineskip}
故
\nopagebreak[4]
\[f\left(\frac{23\pi}6\right)=f\left(\frac{5\pi}6\right)+\sin\frac{5\pi}6 .\]
\Needspace{5\baselineskip}
因为 $0\leqslant\dfrac{5\pi}6<\pi$，所以 $f\left(\dfrac{5\pi}6\right)=0$，于是
\nopagebreak[4]
\[f\left(\frac{23\pi}6\right)=0+\frac12=\frac12 .\]

选 A。

\gkanswer{答案：A。$f(x+3\pi)=f(x)+\sin x$，取 $x=\dfrac{5\pi}6$ 得 $f\left(\dfrac{23\pi}6\right)=\sin\dfrac{5\pi}6=\dfrac12$。}

\end{answer}
\begin{answer}
\gksolhead{7}{}

\textbf{题目}：

一个多面体的三视图如图所示，则该多面体的表面积为\mbox{（　　）}

（A）$21+\sqrt3$；（B）$18+\sqrt3$；（C）$21$；（D）$18$。
\nopagebreak[4]
\begin{center}% ============================================================
% f2014_anhui_science_q07_three_views.tex
% 2014 安徽卷（理科）第 7 题：三视图
% 来源 PDF：普通高考/2014/2014安徽理.pdf
% 原页：第 1 页（A3 横排）中栏上部，印面 857，第 7 题题干上方
% 说明：三个视图都是边长 2 的正方形（边标注 1 与 1）。
%       正视图：实线 (0,1)-(1,0)，虚线 (1,2)-(2,1)；
%       侧视图：虚线 (0,1)-(1,2)，实线 (1,0)-(2,1)；
%       俯视图：实线 (1,2)-(2,1)，虚线 (0,1)-(1,0)。
% ============================================================
\begin{tikzpicture}[>=Stealth, line width=0.7pt, font=\small,
  x=1.55cm, y=1.55cm, line cap=round, line join=round,
  ext/.style={line width=0.4pt},
  dim/.style={<->, line width=0.45pt}]

  % ---------- 正（主）视图 ----------
  \draw[ext] (0,0) -- (2,0) -- (2,2) -- (0,2) -- cycle;
  \draw (0,1) -- (1,0);
  \draw[dashed] (1,2) -- (2,1);
  % 上方 1 + 1
  \draw[ext] (0,2) -- (0,2.55);
  \draw[ext] (1,2) -- (1,2.30);
  \draw[ext] (2,2) -- (2,2.55);
  \draw[dim] (0,2.42) -- (0.92,2.42) node[midway,above,inner sep=1pt] {$1$};
  \draw[dim] (1.08,2.42) -- (2,2.42) node[midway,above,inner sep=1pt] {$1$};
  % 左侧 1 + 1
  \draw[ext] (0,2) -- (-0.55,2);
  \draw[ext] (0,1) -- (-0.30,1);
  \draw[ext] (0,0) -- (-0.55,0);
  \draw[dim] (-0.42,1.08) -- (-0.42,2) node[midway,left,inner sep=1pt] {$1$};
  \draw[dim] (-0.42,0) -- (-0.42,0.92) node[midway,left,inner sep=1pt] {$1$};
  \node at (1,-0.62) {正（主）视图};

  % ---------- 侧（左）视图 ----------
  \begin{scope}[xshift=4.75cm]
    \draw[ext] (0,0) -- (2,0) -- (2,2) -- (0,2) -- cycle;
    \draw[dashed] (0,1) -- (1,2);
    \draw (1,0) -- (2,1);
    % 上方 1 + 1
    \draw[ext] (0,2) -- (0,2.55);
    \draw[ext] (1,2) -- (1,2.30);
    \draw[ext] (2,2) -- (2,2.55);
    \draw[dim] (0,2.42) -- (0.92,2.42) node[midway,above,inner sep=1pt] {$1$};
    \draw[dim] (1.08,2.42) -- (2,2.42) node[midway,above,inner sep=1pt] {$1$};
    % 左侧 1 + 1
    \draw[ext] (0,2) -- (-0.55,2);
    \draw[ext] (0,1) -- (-0.30,1);
    \draw[ext] (0,0) -- (-0.55,0);
    \draw[dim] (-0.42,1.08) -- (-0.42,2) node[midway,left,inner sep=1pt] {$1$};
    \draw[dim] (-0.42,0) -- (-0.42,0.92) node[midway,left,inner sep=1pt] {$1$};
    \node at (1,-0.62) {侧（左）视图};
  \end{scope}

  % ---------- 俯视图 ----------
  \begin{scope}[yshift=-5.15cm]
    \draw[ext] (0,0) -- (2,0) -- (2,2) -- (0,2) -- cycle;
    \draw (1,2) -- (2,1);
    \draw[dashed] (0,1) -- (1,0);
    % 下方 1 + 1
    \draw[ext] (0,0) -- (0,-0.55);
    \draw[ext] (1,0) -- (1,-0.30);
    \draw[ext] (2,0) -- (2,-0.55);
    \draw[dim] (0,-0.42) -- (0.92,-0.42) node[midway,below,inner sep=1pt] {$1$};
    \draw[dim] (1.08,-0.42) -- (2,-0.42) node[midway,below,inner sep=1pt] {$1$};
    % 左侧 1 + 1
    \draw[ext] (0,2) -- (-0.55,2);
    \draw[ext] (0,1) -- (-0.30,1);
    \draw[ext] (0,0) -- (-0.55,0);
    \draw[dim] (-0.42,1.08) -- (-0.42,2) node[midway,left,inner sep=1pt] {$1$};
    \draw[dim] (-0.42,0) -- (-0.42,0.92) node[midway,left,inner sep=1pt] {$1$};
    \node at (1,-0.98) {俯视图};
  \end{scope}

\end{tikzpicture}
\end{center}
\nopagebreak[4]
\Needspace{4\baselineskip}
\textbf{解析}

由三视图还原：该多面体是棱长为 $2$ 的正方体，沿两个相对的顶点各截去一个三条侧棱两两垂直且长均为 $1$ 的四面体。

\Needspace{5\baselineskip}
正方体原表面积为
\nopagebreak[4]
\[6\times2^2=24 .\]
\Needspace{5\baselineskip}
每个截去的四面体去掉三个直角边长为 $1$ 的直角三角形，面积共
\nopagebreak[4]
\[3\times\frac12\times1\times1=\frac32,\]
\Needspace{5\baselineskip}
同时新增一个边长为 $\sqrt2$ 的正三角形，面积为
\nopagebreak[4]
\[\frac{\sqrt3}4(\sqrt2)^2=\frac{\sqrt3}2 .\]
\Needspace{5\baselineskip}
两个四面体共使表面积改变
\nopagebreak[4]
\[2\left(\frac{\sqrt3}2-\frac32\right)=\sqrt3-3 .\]
\Needspace{5\baselineskip}
故表面积为
\nopagebreak[4]
\[24+\sqrt3-3=21+\sqrt3 .\]

选 A。

\gkanswer{答案：A。正方体表面积 $24$，截去两个角上的四面体，各去 $\dfrac32$、增 $\dfrac{\sqrt3}2$，得 $21+\sqrt3$。}

\end{answer}
\begin{answer}
\gksolhead{8}{}

\textbf{题目}：

从正方体六个面的对角线中任取两条作为一对，其中所成的角为 $60^\circ$ 的共有\mbox{（　　）}

（A）$24$ 对；（B）$30$ 对；（C）$48$ 对；（D）$60$ 对。

\Needspace{4\baselineskip}
\textbf{解析}

\Needspace{5\baselineskip}
正方体每个面有 $2$ 条面对角线，共 $12$ 条。按方向分成三族：
\nopagebreak[4]
\[(1,1,0)\ \text{\text{型}},\qquad (1,0,1)\ \text{\text{型}},\qquad (0,1,1)\ \text{\text{型}},\]
每族 $4$ 条，同族内任意两条的夹角为 $0^\circ$ 或 $90^\circ$。

\Needspace{5\baselineskip}
从不同族各取一条，例如取 $(a,b,0)$ 与 $(c,0,d)$，其中 $a,b,c,d\in\{1,-1\}$，则
\nopagebreak[4]
\[(a,b,0)\cdot(c,0,d)=ac=\pm1 .\]
\Needspace{5\baselineskip}
两条对角线的长都是 $\sqrt2$，故夹角 $\alpha$ 满足
\nopagebreak[4]
\[\cos\alpha=\frac{|ac|}{\sqrt2\cdot\sqrt2}=\frac12,\]
即 $\alpha=60^\circ$。因此不同族之间任取两条都符合要求。

\Needspace{5\baselineskip}
不同族的选择有 $\mathrm C_3^2=3$ 种，每族各取一条有 $4\times4=16$ 种，故共有
\nopagebreak[4]
\[3\times16=48\]
对。

选 C。

\gkanswer{答案：C。$12$ 条面对角线分三族，不同族各取一条必成 $60^\circ$，共 $\mathrm C_3^2\times4\times4=48$ 对。}

\end{answer}
\begin{answer}
\gksolhead{9}{}

\textbf{题目}：

若函数 $f(x)=|x+1|+|2x+a|$ 的最小值为 $3$，则实数 $a$ 的值为\mbox{（　　）}

（A）$5$ 或 $8$；（B）$-1$ 或 $5$；（C）$-1$ 或 $-4$；（D）$-4$ 或 $8$。

\Needspace{4\baselineskip}
\textbf{解析}

\Needspace{5\baselineskip}
把函数写成
\nopagebreak[4]
\[f(x)=|x+1|+2\left|x+\frac a2\right| .\]
\Needspace{5\baselineskip}
在 $x<-\dfrac a2$ 时，$2\left|x+\dfrac a2\right|$ 的斜率为 $-2$，$|x+1|$ 的斜率为 $-1$ 或 $1$，故总斜率为 $-3$ 或 $-1$，恒为负；在 $x>-\dfrac a2$ 时，$|x+1|$ 的斜率为 $-1$ 或 $1$，$2\left|x+\dfrac a2\right|$ 的斜率为 $2$，故总斜率为 $1$ 或 $3$，恒为正。所以 $f$ 在 $\left(-\infty,-\dfrac a2\right)$ 上递减、在 $\left(-\dfrac a2,+\infty\right)$ 上递增，最小值在 $x=-\dfrac a2$ 处取得：
\nopagebreak[4]
\[f_{\min}=f\left(-\frac a2\right)=\left|1-\frac a2\right| .\]
\Needspace{5\baselineskip}
由最小值为 $3$ 得
\nopagebreak[4]
\[\left|1-\frac a2\right|=3\iff |2-a|=6,\]
解得 $a=-4$ 或 $a=8$。

选 D。

\gkanswer{答案：D。最小值在 $x=-\dfrac a2$ 取得，$\left|1-\dfrac a2\right|=3$，得 $a=-4$ 或 $8$。}

\end{answer}
\begin{answer}
\gksolhead{10}{}

\textbf{题目}：

在平面直角坐标系 $xOy$ 中，已知向量 $\boldsymbol a$，$\boldsymbol b$，$|\boldsymbol a|=|\boldsymbol b|=1$，$\boldsymbol a\cdot\boldsymbol b=0$，点 $Q$ 满足 $\overrightarrow{OQ}=\sqrt2(\boldsymbol a+\boldsymbol b)$。曲线 $C=\left\{P\mid \overrightarrow{OP}=\boldsymbol a\cos\theta+\boldsymbol b\sin\theta,\ 0\leqslant\theta<2\pi\right\}$，区域 $\Omega=\left\{P\mid 0<r\leqslant|\overrightarrow{PQ}|\leqslant R,\ r<R\right\}$。若 $C\cap\Omega$ 为两段分离的曲线，则\mbox{（　　）}

（A）$1<r<R<3$；（B）$1<r<3\leqslant R$；（C）$r\leqslant1<R<3$；（D）$1<r<3<R$。

\Needspace{4\baselineskip}
\textbf{解析}

\Needspace{5\baselineskip}
由 $\boldsymbol a\perp\boldsymbol b$ 且 $|\boldsymbol a|=|\boldsymbol b|=1$，得
\nopagebreak[4]
\[|\boldsymbol a+\boldsymbol b|^2=1+0+1=2,\]
\Needspace{5\baselineskip}
故
\nopagebreak[4]
\[|OQ|=\sqrt2\cdot\sqrt2=2 .\]
\Needspace{5\baselineskip}
曲线 $C$ 上的点 $P$ 满足 $|OP|=1$，即 $C$ 是以 $O$ 为圆心、半径为 $1$ 的圆。圆心距 $|OQ|=2$，所以圆 $C$ 上的点到 $Q$ 的距离的取值范围是
\nopagebreak[4]
\[[\,2-1,\ 2+1\,]=[1,3].\]

在圆 $C$ 上，条件 $|PQ|\leqslant R$ 截出以“离 $Q$ 最近的点”为中心的一段弧；当 $R<3$ 时它是真弧，当 $R\geqslant3$ 时整个圆都满足。条件 $|PQ|\geqslant r$ 截出以“离 $Q$ 最远的点”为中心的一段弧；当 $r>1$ 时它是真弧，当 $r\leqslant1$ 时整个圆都满足。

\Needspace{5\baselineskip}
两段弧的中点（最近点与最远点）关于 $O$ 对称。要使交集恰为两段分离的曲线，两条边界都必须真正截圆，且交集不能退化为一段，即
\nopagebreak[4]
\[1<r,\qquad R<3 .\]
\Needspace{5\baselineskip}
再结合题设 $r<R$，得
\nopagebreak[4]
\[1<r<R<3 .\]

选 A。

\gkanswer{答案：A。$|OQ|=2$，圆 $C$ 上的点到 $Q$ 的距离在 $[1,3]$ 内，两段分离需 $1<r<R<3$。}

\end{answer}
\begin{answer}
\gksolhead{11}{}

\textbf{题目}：

若将函数 $f(x)=\sin\left(2x+\dfrac\pi4\right)$ 的图象向右平移 $\varphi$ 个单位，所得图象关于 $y$ 轴对称，则 $\varphi$ 的最小正值是\gkblank{}。

\Needspace{4\baselineskip}
\textbf{解析}

\Needspace{5\baselineskip}
设平移后的图象对应的函数为 $g(x)$，则
\nopagebreak[4]
\[g(x)=f(x-\varphi)=\sin\left(2x+\frac\pi4-2\varphi\right).\]
\Needspace{5\baselineskip}
$g(x)$ 的图象关于 $y$ 轴对称，等价于 $g$ 为偶函数。正弦型函数 $\sin(2x+\psi)$ 为偶函数当且仅当
\nopagebreak[4]
\[\psi=\frac\pi2+k\pi,\qquad k\in\mathbf Z .\]
\Needspace{5\baselineskip}
于是
\nopagebreak[4]
\[\frac\pi4-2\varphi=\frac\pi2+k\pi
\iff \varphi=-\frac\pi8-\frac{k\pi}2 .\]
\Needspace{5\baselineskip}
取 $k=-1$ 得最小正值
\nopagebreak[4]
\[\varphi=-\frac\pi8+\frac\pi2=\frac{3\pi}8 .\]

填 $\dfrac{3\pi}8$。

\gkanswer{答案：$\dfrac{3\pi}8$。平移后相位为 $\dfrac\pi4-2\varphi$，令其等于 $\dfrac\pi2+k\pi$，最小正值 $\varphi=\dfrac{3\pi}8$。}

\end{answer}
\begin{answer}
\gksolhead{12}{}

\textbf{题目}：

数列 $\{a_n\}$ 是等差数列，若 $a_1+1$，$a_3+3$，$a_5+5$ 构成公比为 $q$ 的等比数列，则 $q=$\gkblank{}。

\Needspace{4\baselineskip}
\textbf{解析}

\Needspace{5\baselineskip}
设等差数列的公差为 $d$，令 $A=a_1+1$。则
\nopagebreak[4]
\[a_3+3=a_1+2d+3=A+2d+2,\]
\nopagebreak[4]
\[a_5+5=a_1+4d+5=A+4d+4 .\]
\Needspace{5\baselineskip}
三数成等比数列，故
\nopagebreak[4]
\[(A+2d+2)^2=A(A+4d+4).\]
\Needspace{5\baselineskip}
展开左边：
\nopagebreak[4]
\[A^2+4A(d+1)+4(d+1)^2=A^2+4Ad+4A .\]
\Needspace{5\baselineskip}
比较得
\nopagebreak[4]
\[4(d+1)^2=0\iff d=-1 .\]
\Needspace{5\baselineskip}
此时
\nopagebreak[4]
\[a_3+3=A+2(-1)+2=A,\qquad a_5+5=A+4(-1)+4=A,\]
\Needspace{5\baselineskip}
三数相等，故公比
\nopagebreak[4]
\[q=\frac{a_3+3}{a_1+1}=1 .\]

填 $1$。

\gkanswer{答案：$1$。设 $A=a_1+1$，由 $(A+2d+2)^2=A(A+4d+4)$ 得 $d=-1$，三数相等，公比为 $1$。}

\end{answer}
\begin{answer}
\gksolhead{13}{}

\textbf{题目}：

设 $a\neq0$，$n$ 是大于 $1$ 的自然数，$\left(1+\dfrac xa\right)^n$ 的展开式为 $a_0+a_1x+a_2x^2+\cdots+a_nx^n$。若点 $A_i(i,a_i)$（$i=0,1,2$）的位置如图所示，则 $a=$\gkblank{}。
\nopagebreak[4]
\begin{center}% ============================================================
% f2014_anhui_science_q13_point_plot.tex
% 2014 安徽卷（理科）第 13 题：点 A_i(i,a_i) 的坐标图
% 来源 PDF：普通高考/2014/2014安徽理.pdf
% 原页：第 1 页（A3 横排）中栏下部，印面 857，第 13 题题干右侧
% 说明：A_0=(0,1)，A_1=(1,3)，A_2=(2,4)；虚线为向两坐标轴的投影。
% ============================================================
\begin{tikzpicture}[>=Stealth, line width=0.7pt, font=\small,
  x=1.15cm, y=1.15cm, line cap=round, line join=round]

  % 坐标轴
  \draw[->] (-0.45,0) -- (3.35,0) node[right] {$x$};
  \draw[->] (0,-0.45) -- (0,4.85) node[above] {$y$};
  \node[below left] at (0,0) {$O$};

  % 刻度
  \foreach \x in {1,2} {
    \draw (\x,0.06) -- (\x,-0.06);
    \node[below] at (\x,-0.08) {$\x$};
  }
  \foreach \y/\lab in {1/1,3/3,4/4} {
    \draw (0.06,\y) -- (-0.06,\y);
    \node[left] at (-0.08,\y) {$\lab$};
  }

  % 投影虚线
  \draw[dashed] (1,3) -- (1,0);
  \draw[dashed] (2,4) -- (2,0);
  \draw[dashed] (1,3) -- (0,3);
  \draw[dashed] (2,4) -- (0,4);

  % 点
  \fill (0,1) circle (1.6pt) node[right, inner sep=2pt] {$A_0$};
  \fill (1,3) circle (1.6pt) node[right, inner sep=2pt] {$A_1$};
  \fill (2,4) circle (1.6pt) node[right, inner sep=2pt] {$A_2$};

\end{tikzpicture}
\end{center}
\nopagebreak[4]
\Needspace{4\baselineskip}
\textbf{解析}

\Needspace{5\baselineskip}
由二项式定理，
\nopagebreak[4]
\[\left(1+\frac xa\right)^n=\sum_{i=0}^n \mathrm C_n^i\left(\frac xa\right)^i,\]
\Needspace{5\baselineskip}
故
\nopagebreak[4]
\[a_0=1,\qquad a_1=\frac n a,\qquad a_2=\frac{\mathrm C_n^2}{a^2}=\frac{n(n-1)}{2a^2} .\]
\Needspace{5\baselineskip}
图中 $A_0(0,1)$，$A_1(1,3)$，$A_2(2,4)$，即
\nopagebreak[4]
\[a_0=1,\qquad a_1=3,\qquad a_2=4 .\]
\Needspace{5\baselineskip}
由 $a_1=3$ 得 $n=3a$。代入 $a_2=4$：
\nopagebreak[4]
\[\frac{3a(3a-1)}{2a^2}=4
\iff 9a-3=8a
\iff a=3 .\]
此时 $n=9>1$，且 $a_2=\dfrac{9\times8}{2\times9}=4$，与图相符。

填 $3$。

\gkanswer{答案：$3$。$a_0=1$，$a_1=\dfrac na=3$，$a_2=\dfrac{n(n-1)}{2a^2}=4$，联立得 $a=3$。}

\end{answer}
\begin{answer}
\gksolhead{14}{}

\textbf{题目}：

设 $F_1$，$F_2$ 分别是椭圆 $E:x^2+\dfrac{y^2}{b^2}=1$（$0<b<1$）的左、右焦点，过点 $F_1$ 的直线交椭圆 $E$ 于 $A$，$B$ 两点，若 $|AF_1|=3|BF_1|$，$AF_2\perp x$ 轴，则椭圆 $E$ 的方程为\gkblank{}。

\Needspace{4\baselineskip}
\textbf{解析}

\Needspace{5\baselineskip}
由椭圆方程 $x^2+\dfrac{y^2}{b^2}=1$ 知长半轴为 $1$，故
\nopagebreak[4]
\[c^2=1-b^2,\qquad F_1(-c,0),\quad F_2(c,0).\]
\Needspace{5\baselineskip}
因为 $AF_2\perp x$ 轴，所以 $A$ 的横坐标为 $c$。代入椭圆方程：
\nopagebreak[4]
\[c^2+\frac{y^2}{b^2}=1\Longrightarrow y^2=b^2(1-c^2)=b^4,\]
不妨取 $A(c,b^2)$。

\Needspace{5\baselineskip}
直线 $F_1A$ 的参数方程为
\nopagebreak[4]
\[(x,y)=(-c,0)+t(2c,b^2),\]
\Needspace{5\baselineskip}
其中 $t=1$ 对应点 $A$，$t=0$ 对应 $F_1$。由于 $F_1$ 在椭圆内部，直线与椭圆的另一交点 $B$ 对应 $t<0$。由 $|AF_1|=3|BF_1|$ 得
\nopagebreak[4]
\[|t_B|=\frac13,\qquad t_B=-\frac13 .\]

\Needspace{5\baselineskip}
把参数式代入椭圆方程：
\nopagebreak[4]
\[c^2(2t-1)^2+b^2t^2=1,\]
\Needspace{5\baselineskip}
整理得
\nopagebreak[4]
\[(4c^2+b^2)t^2-4c^2t-b^2=0 .\]
\Needspace{5\baselineskip}
它的两个根为 $1$ 与 $-\dfrac13$，由韦达定理
\nopagebreak[4]
\[\frac{-b^2}{4c^2+b^2}=-\frac13
\iff 3b^2=4c^2+b^2
\iff b^2=2c^2 .\]
\Needspace{5\baselineskip}
又 $c^2=1-b^2$，所以
\nopagebreak[4]
\[b^2=2(1-b^2)\iff 3b^2=2\iff b^2=\frac23 .\]
\Needspace{5\baselineskip}
故椭圆方程为
\nopagebreak[4]
\[x^2+\frac{3y^2}{2}=1 .\]

填 $x^2+\dfrac{3y^2}{2}=1$。

\gkanswer{答案：$x^2+\dfrac{3y^2}{2}=1$。取 $A(c,b^2)$，直线参数代入椭圆，由两根 $1$ 与 $-\dfrac13$ 得 $b^2=2c^2$，结合 $c^2=1-b^2$ 得 $b^2=\dfrac23$。}

\end{answer}
\begin{answer}
\gksolhead{15}{}

\textbf{题目}：

已知两个不相等的非零向量 $\boldsymbol a$，$\boldsymbol b$，两组向量 $\boldsymbol x_1$，$\boldsymbol x_2$，$\boldsymbol x_3$，$\boldsymbol x_4$，$\boldsymbol x_5$ 和 $\boldsymbol y_1$，$\boldsymbol y_2$，$\boldsymbol y_3$，$\boldsymbol y_4$，$\boldsymbol y_5$ 均由 $2$ 个 $\boldsymbol a$ 和 $3$ 个 $\boldsymbol b$ 排列而成，记 $S=\boldsymbol x_1\cdot\boldsymbol y_1+\boldsymbol x_2\cdot\boldsymbol y_2+\boldsymbol x_3\cdot\boldsymbol y_3+\boldsymbol x_4\cdot\boldsymbol y_4+\boldsymbol x_5\cdot\boldsymbol y_5$，$S_{\min}$ 表示 $S$ 所有可能取值中的最小值。则下列命题正确的是\gkblank{}。（写出所有正确命题的编号）

① $S$ 有 $5$ 个不同的值；

② 若 $\boldsymbol a\perp\boldsymbol b$，则 $S_{\min}$ 与 $|\boldsymbol a|$ 无关；

③ 若 $\boldsymbol a\parallel\boldsymbol b$，则 $S_{\min}$ 与 $|\boldsymbol b|$ 无关；

④ 若 $|\boldsymbol b|>4|\boldsymbol a|$，则 $S_{\min}>0$；

⑤ 若 $|\boldsymbol b|=2|\boldsymbol a|$，$S_{\min}=8|\boldsymbol a|^2$，则 $\boldsymbol a$ 与 $\boldsymbol b$ 的夹角为 $\dfrac\pi4$。

\Needspace{4\baselineskip}
\textbf{解析}

\Needspace{5\baselineskip}
记 $A=|\boldsymbol a|^2$，$B=|\boldsymbol b|^2$，$C=\boldsymbol a\cdot\boldsymbol b$。设 $k$ 为 $\boldsymbol x_i=\boldsymbol y_i=\boldsymbol a$ 的位置个数。由两组向量各含 $2$ 个 $\boldsymbol a$、$3$ 个 $\boldsymbol b$，得各类位置数为
\nopagebreak[4]
\[(\boldsymbol a,\boldsymbol a):k,\quad (\boldsymbol a,\boldsymbol b):2-k,\quad
(\boldsymbol b,\boldsymbol a):2-k,\quad (\boldsymbol b,\boldsymbol b):1+k .\]
\Needspace{5\baselineskip}
因此
\nopagebreak[4]
\[S=kA+(4-2k)C+(1+k)B=k(A-2C+B)+B+4C,\]
其中 $k=0,1,2$。

① 错：$S$ 至多有 $3$ 个不同取值，不可能有 $5$ 个。

\Needspace{5\baselineskip}
② 对：若 $\boldsymbol a\perp\boldsymbol b$，则 $C=0$，
\nopagebreak[4]
\[S=kA+B,\]
当 $k=0$ 时取最小值 $S_{\min}=B=|\boldsymbol b|^2$，与 $|\boldsymbol a|$ 无关。

\Needspace{5\baselineskip}
③ 错：若 $\boldsymbol a\parallel\boldsymbol b$，设 $\boldsymbol b=\lambda\boldsymbol a$（$\lambda\neq1$），则 $C=\lambda A$，$B=\lambda^2A$，
\nopagebreak[4]
\[S=A\bigl[\lambda^2+4\lambda+k(\lambda-1)^2\bigr].\]
\Needspace{5\baselineskip}
因为 $(\lambda-1)^2>0$，最小值在 $k=0$ 取得：
\nopagebreak[4]
\[S_{\min}=A\lambda^2+4A\lambda=B+4\lambda A .\]
它随 $|\boldsymbol b|$ 改变，故与 $|\boldsymbol b|$ 有关。

\Needspace{5\baselineskip}
④ 对：由 $\boldsymbol a\neq\boldsymbol b$，
\nopagebreak[4]
\[A-2C+B=|\boldsymbol a-\boldsymbol b|^2>0,\]
\Needspace{5\baselineskip}
所以 $S$ 随 $k$ 增大而增大，最小值在 $k=0$：
\nopagebreak[4]
\[S_{\min}=B+4C\geqslant B-4|\boldsymbol a||\boldsymbol b|
=|\boldsymbol b|\bigl(|\boldsymbol b|-4|\boldsymbol a|\bigr)>0 .\]

\Needspace{5\baselineskip}
⑤ 错：若 $|\boldsymbol b|=2|\boldsymbol a|$，则 $B=4A$，且
\nopagebreak[4]
\[S_{\min}=B+4C=4A+4C .\]
\Needspace{5\baselineskip}
由 $S_{\min}=8A$ 得 $C=A$，于是
\nopagebreak[4]
\[\cos\theta=\frac{C}{|\boldsymbol a||\boldsymbol b|}
=\frac{A}{2A}=\frac12,\]
夹角应为 $\dfrac\pi3$，不是 $\dfrac\pi4$。

故正确的是②④。

\gkanswer{答案：②④。$S=k|\boldsymbol a-\boldsymbol b|^2+|\boldsymbol b|^2+4\boldsymbol a\cdot\boldsymbol b$，$k=0,1,2$；②中 $S_{\min}=|\boldsymbol b|^2$，④由 $|\boldsymbol b|(|\boldsymbol b|-4|\boldsymbol a|)>0$ 得证。}

\end{answer}
\begin{answer}
\gksolhead{16}{}

\textbf{题目}：

设 $\triangle ABC$ 的内角 $A$，$B$，$C$ 所对边的长分别是 $a$，$b$，$c$，且 $b=3$，$c=1$，$A=2B$。

（1） 求 $a$ 的值；

（2） 求 $\sin\left(A+\dfrac\pi4\right)$ 的值。

\Needspace{4\baselineskip}
\textbf{解析}

\Needspace{5\baselineskip}
\textbf{（1）} 由正弦定理
\nopagebreak[4]
\[\frac a{\sin A}=\frac b{\sin B},\qquad A=2B,\quad b=3,\]
\Needspace{5\baselineskip}
得
\nopagebreak[4]
\[a=\frac{3\sin2B}{\sin B}=6\cos B .\]
\Needspace{5\baselineskip}
再由余弦定理
\nopagebreak[4]
\[b^2=a^2+c^2-2ac\cos B,\]
\Needspace{5\baselineskip}
代入 $b=3$，$c=1$，$\cos B=\dfrac a6$：
\nopagebreak[4]
\[9=a^2+1-2a\cdot\frac a6
=a^2+1-\frac{a^2}3
=\frac23a^2+1 .\]
\Needspace{5\baselineskip}
所以
\nopagebreak[4]
\[\frac23a^2=8\Longrightarrow a^2=12\Longrightarrow a=2\sqrt3 .\]
此时 $\cos B=\dfrac{a}{6}=\dfrac{\sqrt3}{3}$，符合三角形边角关系。

\Needspace{5\baselineskip}
\textbf{（2）} 由余弦定理
\nopagebreak[4]
\[\cos A=\frac{b^2+c^2-a^2}{2bc}
=\frac{9+1-12}{2\times3\times1}
=-\frac13 .\]
\Needspace{5\baselineskip}
因为 $A\in(0,\pi)$，所以
\nopagebreak[4]
\[\sin A=\sqrt{1-\cos^2A}=\sqrt{1-\frac19}=\frac{2\sqrt2}{3} .\]
\Needspace{5\baselineskip}
于是
\nopagebreak[4]
\[\sin\left(A+\frac\pi4\right)
=\sin A\cos\frac\pi4+\cos A\sin\frac\pi4
=\frac{\sin A+\cos A}{\sqrt2}
=\frac{\dfrac{2\sqrt2}{3}-\dfrac13}{\sqrt2}
=\frac{4-\sqrt2}{6} .\]

\gkanswer{答案（1）：$a=2\sqrt3$。由 $a=6\cos B$ 与余弦定理 $9=a^2+1-\dfrac{a^2}3$ 解得。}

\gkanswer{答案（2）：$\dfrac{4-\sqrt2}{6}$。$\cos A=-\dfrac13$，$\sin A=\dfrac{2\sqrt2}3$，再用和角公式。}

\end{answer}
\begin{answer}
\gksolhead{17}{}

\textbf{题目}：

甲乙两人进行围棋比赛，约定先连胜两局者直接赢得比赛，若赛完 $5$ 局仍未出现连胜，则判定获胜局数多者赢得比赛。假设每局甲获胜的概率为 $\dfrac23$，乙获胜的概率为 $\dfrac13$，各局比赛结果相互独立。

（1） 求甲在 $4$ 局以内（含 $4$ 局）赢得比赛的概率；

（2） 记 $X$ 为比赛决出胜负时的总局数，求 $X$ 的分布列和均值（数学期望）。

\Needspace{4\baselineskip}
\textbf{解析}

记甲胜一局的概率为 $p=\dfrac23$，乙胜一局的概率为 $q=\dfrac13$，且 $p+q=1$。

\textbf{（1）} 甲在 $4$ 局以内获胜，分三种情形：

甲2局获胜：前两局甲连胜，概率 $p^2$；

甲3局获胜：只能是“乙、甲、甲”，概率 $qp^2$；

甲4局获胜：第3、4局必须都是甲胜。若第2局为甲胜，则第3局不能再为甲胜（否则甲在第3局已连胜），而第3局若为乙胜又无法在第4局形成甲的连胜，故第2局只能为乙胜；此时若第1局为乙胜，则乙已在第2局连胜，故第1局必须为甲胜。唯一序列为“甲、乙、甲、甲”，概率为 $pqp^2$。

\Needspace{5\baselineskip}
因此
\nopagebreak[4]
\[P=p^2+qp^2+pqp^2
=p^2(1+q+pq)
=\frac49\left(1+\frac13+\frac29\right)
=\frac49\cdot\frac{14}{9}
=\frac{56}{81} .\]

\textbf{（2）} $X$ 的可能取值为 $2,3,4,5$。

\Needspace{5\baselineskip}
$X=2$：前两局同一人连胜，
\nopagebreak[4]
\[P(X=2)=p^2+q^2=\frac49+\frac19=\frac59 .\]

\Needspace{5\baselineskip}
$X=3$：前两局不同，第3局与第2局相同，即“乙甲甲”或“甲乙乙”，
\nopagebreak[4]
\[P(X=3)=qp^2+pq^2=pq(p+q)=pq=\frac29 .\]

\Needspace{5\baselineskip}
$X=4$：前3局不能出现连胜，第4局与第3局相同。甲的取胜序列只能是“甲乙甲甲”，乙的取胜序列只能是“乙甲乙乙”，故
\nopagebreak[4]
\[P(X=4)=p^3q+q^3p=pq(p^2+q^2)
=\frac29\cdot\frac59=\frac{10}{81} .\]

\Needspace{5\baselineskip}
$X=5$：包括两类。第一类是前4局交替、第5局与第4局相同：“甲乙甲乙乙”或“乙甲乙甲甲”，概率
\nopagebreak[4]
\[p^2q^3+q^2p^3=p^2q^2(p+q)=\frac49\cdot\frac19=\frac4{81};\]
\Needspace{5\baselineskip}
第二类是5局全部交替，未出现连胜，由胜局多者获胜：“甲乙甲乙甲”（甲胜3局）或“乙甲乙甲乙”（乙胜3局），概率同为
\nopagebreak[4]
\[p^3q^2+q^3p^2=p^2q^2(p+q)=\frac4{81}.\]
\Needspace{5\baselineskip}
故
\nopagebreak[4]
\[P(X=5)=\frac4{81}+\frac4{81}=\frac8{81} .\]

分布列为

\begin{center}
\begin{tabular}{lllll}
\toprule
$X$ & $2$ & $3$ & $4$ & $5$ \\
\midrule
$P$ & $\dfrac59$ & $\dfrac29$ & $\dfrac{10}{81}$ & $\dfrac8{81}$ \\
\bottomrule
\end{tabular}
\end{center}

\Needspace{5\baselineskip}
均值为
\nopagebreak[4]
\[E(X)=2\cdot\frac59+3\cdot\frac29+4\cdot\frac{10}{81}+5\cdot\frac8{81}
=\frac{90+54+40+40}{81}
=\frac{224}{81} .\]

\gkanswer{答案（1）：$\dfrac{56}{81}$。甲2局、3局、4局获胜的概率分别为 $\dfrac49$、$\dfrac4{27}$、$\dfrac8{81}$，相加。}

\gkanswer{答案（2）：$P(X=2)=\dfrac59$，$P(X=3)=\dfrac29$，$P(X=4)=\dfrac{10}{81}$，$P(X=5)=\dfrac8{81}$；$E(X)=\dfrac{224}{81}$。按首次出现连胜的序列计数。}

\end{answer}
\begin{answer}
\gksolhead{18}{}

\textbf{题目}：

设函数 $f(x)=1+(1+a)x-x^2-x^3$，其中 $a>0$。

（1） 讨论 $f(x)$ 在其定义域上的单调性；

（2） 当 $x\in[0,1]$ 时，求 $f(x)$ 取得最大值和最小值时的 $x$ 的值。

\Needspace{4\baselineskip}
\textbf{解析}

\Needspace{5\baselineskip}
\textbf{（1）} 定义域为 $\mathbf R$，且
\nopagebreak[4]
\[f'(x)=1+a-2x-3x^2=-(3x^2+2x-1-a).\]
\Needspace{5\baselineskip}
方程 $3x^2+2x-(1+a)=0$ 的两根为
\nopagebreak[4]
\[x_1=\frac{-1-\sqrt{4+3a}}3,\qquad
x_2=\frac{-1+\sqrt{4+3a}}3 .\]
\Needspace{5\baselineskip}
因为 $a>0$，所以 $\sqrt{4+3a}>2$，从而
\nopagebreak[4]
\[x_1<-1,\qquad x_2>\frac13 .\]
二次项系数为负，故 $f'(x)<0$ 当 $x<x_1$ 或 $x>x_2$，$f'(x)>0$ 当 $x_1<x<x_2$。于是 $f(x)$ 在 $(-\infty,x_1)$ 上单调递减，在 $(x_1,x_2)$ 上单调递增，在 $(x_2,+\infty)$ 上单调递减。

\Needspace{5\baselineskip}
\textbf{（2）} 由 $x_2=\dfrac{\sqrt{4+3a}-1}{3}$ 得
\nopagebreak[4]
\[x_2<1\iff \sqrt{4+3a}<4\iff a<4 .\]
又 $f(0)=1$，$f(1)=1+(1+a)-1-1=a$。

当 $0<a<4$ 时，$x_2\in(0,1)$，$f$ 在 $[0,x_2]$ 上递增、在 $[x_2,1]$ 上递减，故最大值在 $x=x_2$ 取得。最小值在端点取得：若 $0<a<1$，则 $f(1)=a<1=f(0)$，最小值在 $x=1$；若 $a=1$，两端点函数值都为 $1$，最小值在 $x=0$ 或 $x=1$；若 $1<a<4$，则 $f(0)=1<a=f(1)$，最小值在 $x=0$。

当 $a\geqslant4$ 时，$x_2\geqslant1$，$f$ 在 $[0,1]$ 上单调递增，故最大值在 $x=1$ 取得，最小值在 $x=0$ 取得。

综上：最大值点当 $0<a<4$ 时为 $x=\dfrac{\sqrt{4+3a}-1}{3}$，当 $a\geqslant4$ 时为 $x=1$；最小值点当 $0<a<1$ 时为 $x=1$，当 $a=1$ 时为 $x=0$ 或 $x=1$，当 $a>1$ 时为 $x=0$。

\gkanswer{答案（1）：记 $x_1=\dfrac{-1-\sqrt{4+3a}}3$，$x_2=\dfrac{-1+\sqrt{4+3a}}3$；减区间为 $(-\infty,x_1)$、$(x_2,+\infty)$，增区间为 $(x_1,x_2)$。由 $f'(x)=-(3x^2+2x-1-a)$ 的符号得。}

\gkanswer{答案（2）：最大值点当 $0<a<4$ 时为 $x=\dfrac{\sqrt{4+3a}-1}{3}$，当 $a\geqslant4$ 时为 $x=1$；最小值点当 $0<a<1$ 时为 $x=1$，当 $a=1$ 时为 $x=0$ 或 $x=1$，当 $a>1$ 时为 $x=0$。比较 $f(0)=1$、$f(1)=a$ 与极大值点。}

\end{answer}
\begin{answer}
\gksolhead{19}{}

\textbf{题目}：

如图，已知两条抛物线 $E_1:y^2=2p_1x$（$p_1>0$）和 $E_2:y^2=2p_2x$（$p_2>0$），过原点 $O$ 的两条直线 $l_1$ 和 $l_2$，$l_1$ 与 $E_1$，$E_2$ 分别交于 $A_1$，$A_2$ 两点，$l_2$ 与 $E_1$，$E_2$ 分别交于 $B_1$，$B_2$ 两点。
\nopagebreak[4]
\begin{center}% ============================================================
% f2014_anhui_science_q19_parabolas.tex
% 2014 安徽卷（理科）第 19 题：两条抛物线
% 来源 PDF：普通高考/2014/2014安徽理.pdf
% 原页：第 2 页（A3 横排）左栏下部，印面 858，第 19 题题干下方
% 说明：E_1: y^2=2p_1x，E_2: y^2=2p_2x（图中 p_2>p_1）。
%       l_1 与 E_1,E_2 交于 A_1,A_2；l_2 与 E_1,E_2 交于 B_1,B_2。
%       图中另画出 A_1B_1 与 A_2B_2 两条弦。
% ============================================================
\begin{tikzpicture}[>=Stealth, line width=0.7pt, font=\small,
  x=1.45cm, y=1.10cm, line cap=round, line join=round]

  % 坐标轴
  \draw[->] (-0.45,0) -- (3.15,0) node[right] {$x$};
  \draw[->] (0,-3.55) -- (0,3.75) node[above] {$y$};
  \node[below left] at (0,0) {$O$};

  % 抛物线（上、下两支）
  \draw[domain=0:2.55, samples=120, smooth, variable=\t]
    plot ({\t},{sqrt(2*\t)});
  \draw[domain=0:2.55, samples=120, smooth, variable=\t]
    plot ({\t},{-sqrt(2*\t)});
  \draw[domain=0:2.55, samples=120, smooth, variable=\t]
    plot ({\t},{sqrt(4*\t)});
  \draw[domain=0:2.55, samples=120, smooth, variable=\t]
    plot ({\t},{-sqrt(4*\t)});

  % 两条过原点的直线
  \draw (-0.28,-0.56) -- (1.05,2.10);
  \draw (-0.32,0.48) -- (2.05,-3.08);

  % A1B1、A2B2
  \draw[line width=0.5pt] (0.5,1) -- (0.8889,-1.3333);
  \draw[line width=0.5pt] (1,2) -- (1.7778,-2.6667);

  % 交点
  \fill (0.5,1) circle (1.4pt) node[below right, inner sep=2pt] {$A_1$};
  \fill (1,2) circle (1.4pt) node[above left, inner sep=3pt] {$A_2$};
  \fill (0.8889,-1.3333) circle (1.4pt) node at (1.10,-1.22) {$B_1$};
  \fill (1.7778,-2.6667) circle (1.4pt) node[below left, inner sep=3pt] {$B_2$};

  % 曲线与直线标签
  \node[right] at (2.55,2.23) {$E_1$};
  \node[right] at (2.55,3.16) {$E_2$};
  \node[left] at (-0.34,-0.70) {$l_1$};
  \node[left] at (-0.40,0.62) {$l_2$};

\end{tikzpicture}
\end{center}
\nopagebreak[4]
（1） 证明：$A_1B_1\parallel A_2B_2$；

（2） 过原点 $O$ 作直线 $l$（异于 $l_1$，$l_2$）与 $E_1$，$E_2$ 分别交于 $C_1$，$C_2$ 两点。记 $\triangle A_1B_1C_1$ 与 $\triangle A_2B_2C_2$ 的面积分别为 $S_1$ 与 $S_2$，求 $\dfrac{S_1}{S_2}$ 的值。

\Needspace{4\baselineskip}
\textbf{解析}

\Needspace{5\baselineskip}
\textbf{（1）} 设直线 $l_1$，$l_2$ 的斜率分别为 $k_1$，$k_2$，则
\nopagebreak[4]
\[l_1:y=k_1x,\qquad l_2:y=k_2x .\]
\Needspace{5\baselineskip}
它们与抛物线 $y^2=2px$ 的交点满足
\nopagebreak[4]
\[k^2x^2=2px\Longrightarrow x=0\ \text{\text{或}}\ x=\frac{2p}{k^2},\]
\Needspace{5\baselineskip}
其中 $x=0$ 对应原点 $O$。故非原点的交点为
\nopagebreak[4]
\[\left(\frac{2p}{k^2},\frac{2p}{k}\right).\]
\Needspace{5\baselineskip}
于是
\nopagebreak[4]
\[A_i=\left(\frac{2p_i}{k_1^2},\frac{2p_i}{k_1}\right),\qquad
B_i=\left(\frac{2p_i}{k_2^2},\frac{2p_i}{k_2}\right)\quad (i=1,2).\]
\Needspace{5\baselineskip}
从而
\nopagebreak[4]
\[\overrightarrow{A_iB_i}
=2p_i\left(\frac1{k_2^2}-\frac1{k_1^2},\ \frac1{k_2}-\frac1{k_1}\right).\]
\Needspace{5\baselineskip}
所以
\nopagebreak[4]
\[\overrightarrow{A_1B_1}=\frac{p_1}{p_2}\overrightarrow{A_2B_2},\]
即 $A_1B_1\parallel A_2B_2$。

\Needspace{5\baselineskip}
\textbf{（2）} 设直线 $l$ 的斜率为 $k$，则
\nopagebreak[4]
\[C_i=\left(\frac{2p_i}{k^2},\frac{2p_i}{k}\right)\quad (i=1,2).\]
\Needspace{5\baselineskip}
比较 $i=1,2$ 的坐标可见，把 $E_1$ 上的三个点 $A_1$，$B_1$，$C_1$ 的坐标统一乘以 $\dfrac{p_2}{p_1}$，就得到 $E_2$ 上的三个点 $A_2$，$B_2$，$C_2$ 的坐标。以原点为中心的位似变换把三角形 $A_1B_1C_1$ 变为三角形 $A_2B_2C_2$，位似比为 $\dfrac{p_2}{p_1}$，故面积比为位似比的平方：
\nopagebreak[4]
\[\frac{S_1}{S_2}=\left(\frac{p_1}{p_2}\right)^2 .\]

\gkanswer{答案（1）：见解析版}

\gkanswer{答案（2）：$\dfrac{p_1^2}{p_2^2}$。三顶点坐标都随 $p$ 同比缩放，面积比为相似比的平方。}

\end{answer}
\begin{answer}
\gksolhead{20}{}

\textbf{题目}：

如图，四棱柱 $ABCD-A_1B_1C_1D_1$ 中，$A_1A\perp$ 底面 $ABCD$。四边形 $ABCD$ 为梯形，$AD\parallel BC$，且 $AD=2BC$。过 $A_1$，$C$，$D$ 三点的平面记为 $\alpha$，$BB_1$ 与 $\alpha$ 的交点为 $Q$。
\nopagebreak[4]
\begin{center}% ============================================================
% f2014_anhui_science_q20_prism.tex
% 2014 安徽卷（理科）第 20 题：四棱柱与截面
% 来源 PDF：普通高考/2014/2014安徽理.pdf
% 原页：第 2 页（A3 横排）中栏中部，印面 858，第 20 题题干下方
% 说明：四棱柱 ABCD-A_1B_1C_1D_1，A_1A⊥底面，AD∥BC 且 AD=2BC。
%       平面 α 过 A_1,C,D，交 BB_1 于 Q；截面为四边形 A_1QCD。
%       被遮挡的 AD、DA_1 画虚线，其余轮廓与 A_1Q、QC、CD 画实线。
% ============================================================
\begin{tikzpicture}[>=Stealth, line width=0.7pt, font=\small,
  x=1.25cm, y=1.25cm, line cap=round, line join=round]

  % 顶点
  \coordinate (A)  at (0,0);
  \coordinate (B)  at (0.35,-0.35);
  \coordinate (C)  at (1.55,-0.35);
  \coordinate (D)  at (2.40,0);
  \coordinate (A1) at (0,2.40);
  \coordinate (B1) at (0.35,2.05);
  \coordinate (C1) at (1.55,2.05);
  \coordinate (D1) at (2.40,2.40);
  \coordinate (Q)  at (0.35,1.025);

  % 截面淡色填充
  \fill[gray!12] (A1) -- (Q) -- (C) -- (D) -- cycle;

  % 被遮挡的棱与截面边
  \draw[dashed] (A) -- (D);
  \draw[dashed] (D) -- (A1);

  % 底面
  \draw (A) -- (B) -- (C) -- (D);
  % 四条侧棱
  \draw (A) -- (A1);
  \draw (B) -- (B1);
  \draw (C) -- (C1);
  \draw (D) -- (D1);
  % 上底面
  \draw (A1) -- (B1) -- (C1) -- (D1) -- cycle;
  % 截面在可见面上的边
  \draw (A1) -- (Q) -- (C);

  % 顶点与截面标签
  \node[left]  at (A)  {$A$};
  \node[below] at (B)  {$B$};
  \node[below] at (C)  {$C$};
  \node[right] at (D)  {$D$};
  \node[left]  at (A1) {$A_1$};
  \node[above right, inner sep=2pt] at (B1) {$B_1$};
  \node[below left, inner sep=2pt] at (C1) {$C_1$};
  \node[right] at (D1) {$D_1$};
  \node[left=2pt] at (Q)  {$Q$};
  \node at (1.12,0.78) {$\alpha$};

\end{tikzpicture}
\end{center}
\nopagebreak[4]
（1） 证明：$Q$ 为 $BB_1$ 的中点；

（2） 求此四棱柱被平面 $\alpha$ 所分成上下两部分的体积之比；

（3） 若 $A_1A=4$，$CD=2$，梯形 $ABCD$ 的面积为 $6$，求平面 $\alpha$ 与底面 $ABCD$ 所成二面角的大小。

\Needspace{4\baselineskip}
\textbf{解析}

\Needspace{5\baselineskip}
\textbf{（1）} 建立空间直角坐标系。令
\nopagebreak[4]
\[A(0,0,0),\quad D(2,0,0),\quad B(u,v,0),\quad C(u+1,v,0),\]
其中 $v\neq0$，并设 $A_1(0,0,h)$。这里取 $AD=2$、$BC=1$，不影响结论。

\Needspace{5\baselineskip}
平面 $\alpha$ 过 $A_1$，$C$，$D$，其中
\nopagebreak[4]
\[\overrightarrow{A_1C}=(u+1,v,-h),\qquad
\overrightarrow{A_1D}=(2,0,-h).\]
\Needspace{5\baselineskip}
取法向量
\nopagebreak[4]
\[\boldsymbol n=\overrightarrow{A_1C}\times\overrightarrow{A_1D}
=(-hv,\ h(u-1),\ -2v).\]
\Needspace{5\baselineskip}
故 $\alpha$ 的方程为
\nopagebreak[4]
\[-hvx+h(u-1)y-2v(z-h)=0 .\]
\Needspace{5\baselineskip}
点 $Q$ 在 $BB_1$ 上，可设 $Q(u,v,t)$。代入平面方程：
\nopagebreak[4]
\[-huv+h(u-1)v-2vt+2vh=0,\]
\Needspace{5\baselineskip}
即
\nopagebreak[4]
\[-hv-2vt+2vh=0\Longrightarrow t=\frac h2 .\]
所以 $Q$ 的竖坐标为 $\dfrac h2$，即 $Q$ 为 $BB_1$ 的中点。

\Needspace{5\baselineskip}
\textbf{（2）} 体积比在可逆仿射变换下不变，故可取
\nopagebreak[4]
\[A(0,0,0),\quad B(0,1,0),\quad C(1,1,0),\quad D(2,0,0),\]
\Needspace{5\baselineskip}
并设棱柱高为 $H$。此时 $A_1(0,0,H)$，平面 $\alpha$ 的方程为
\nopagebreak[4]
\[Hx+Hy+2z=2H .\]
\Needspace{5\baselineskip}
含底面 $ABCD$ 的下半部分满足
\nopagebreak[4]
\[z\leqslant\frac H2(2-x-y).\]
\Needspace{5\baselineskip}
底面梯形 $ABCD$ 的面积为
\nopagebreak[4]
\[S=\frac{AD+BC}2\times1=\frac{2+1}2=\frac32 .\]
\Needspace{5\baselineskip}
直接计算
\nopagebreak[4]
\[\iint_{ABCD}x\,\mathrm dA=\int_0^1\int_0^{2-y}x\,\mathrm dx\,\mathrm dy
=\int_0^1\frac{(2-y)^2}2\,\mathrm dy=\frac76,\]
\nopagebreak[4]
\[\iint_{ABCD}y\,\mathrm dA=\int_0^1y(2-y)\,\mathrm dy=\frac23 .\]
\Needspace{5\baselineskip}
于是下半部分体积为
\nopagebreak[4]
\[V_{\text{\text{下}}}=\frac H2\iint_{ABCD}(2-x-y)\,\mathrm dA
=\frac H2\left(2\cdot\frac32-\frac76-\frac23\right)
=\frac H2\cdot\frac76=\frac{7H}{12}.\]
\Needspace{5\baselineskip}
整个棱柱的体积为
\nopagebreak[4]
\[V=\frac32H=\frac{18H}{12},\]
\Needspace{5\baselineskip}
故上半部分体积为 $\dfrac{11H}{12}$。所以上、下两部分的体积之比为
\nopagebreak[4]
\[V_{\text{\text{上}}}:V_{\text{\text{下}}}=11:7 .\]

\Needspace{5\baselineskip}
\textbf{（3）} 设平面 $\alpha$ 与底面所成的二面角为 $\theta$。因为 $A_1A\perp$ 底面，$A$ 是 $A_1$ 在底面上的射影，所以
\nopagebreak[4]
\[\tan\theta=\frac{A_1A}{d(A,CD)},\]
其中 $d(A,CD)$ 是点 $A$ 到直线 $CD$ 的距离。

\Needspace{5\baselineskip}
梯形 $ABCD$ 的面积为 $6$。对角线 $AC$ 把它分成 $\triangle ABC$ 与 $\triangle ACD$，两者高相同，底的比为
\nopagebreak[4]
\[BC:AD=1:2,\]
\Needspace{5\baselineskip}
故
\nopagebreak[4]
\[S_{\triangle ACD}=\frac23\times6=4 .\]
\Needspace{5\baselineskip}
另一方面，
\nopagebreak[4]
\[S_{\triangle ACD}=\frac12\cdot CD\cdot d(A,CD)
=\frac12\times2\times d(A,CD)=d(A,CD),\]
\Needspace{5\baselineskip}
所以 $d(A,CD)=4$。又 $A_1A=4$，于是
\nopagebreak[4]
\[\tan\theta=\frac44=1\Longrightarrow\theta=\frac\pi4 .\]
故平面 $\alpha$ 与底面 $ABCD$ 所成二面角为 $\dfrac\pi4$。

\gkanswer{答案（1）：见解析版}

\gkanswer{答案（2）：上、下两部分的体积之比为 $11:7$。取标准梯形积分，下半体积为 $\dfrac{7H}{12}$，上半为 $\dfrac{11H}{12}$。}

\gkanswer{答案（3）：$\dfrac\pi4$。$S_{\triangle ACD}=4=\dfrac12\cdot CD\cdot d(A,CD)$ 得 $d(A,CD)=4$，故 $\tan\theta=\dfrac{A_1A}{d}=1$。}

\end{answer}
\begin{answer}
\gksolhead{21}{}

\textbf{题目}：

设实数 $c>0$，整数 $p>1$，$n\in\mathbf N^*$。

（1） 证明：当 $x>-1$ 且 $x\neq0$ 时，$(1+x)^p>1+px$；

（2） 数列 $\{a_n\}$ 满足 $a_1>c^{\frac1p}$，$a_{n+1}=\dfrac{p-1}p a_n+\dfrac cp a_n^{1-p}$，证明：$a_n>a_{n+1}>c^{\frac1p}$。

\Needspace{4\baselineskip}
\textbf{解析}

\Needspace{5\baselineskip}
\textbf{（1）} 设
\nopagebreak[4]
\[f(x)=(1+x)^p-1-px,\qquad x>-1 .\]
\Needspace{5\baselineskip}
则 $f(0)=0$，且
\nopagebreak[4]
\[f'(x)=p(1+x)^{p-1}-p=p\left[(1+x)^{p-1}-1\right].\]
因为 $p>1$，所以 $p-1>0$。

\Needspace{5\baselineskip}
当 $-1<x<0$ 时，$0<1+x<1$，故 $(1+x)^{p-1}<1$，从而 $f'(x)<0$；当 $x>0$ 时，$1+x>1$，故 $(1+x)^{p-1}>1$，从而 $f'(x)>0$。因此 $f$ 在 $(-1,0)$ 上递减、在 $(0,+\infty)$ 上递增，最小值在 $x=0$ 取得：
\nopagebreak[4]
\[f(x)>f(0)=0\qquad (x>-1,\ x\neq0).\]
\Needspace{5\baselineskip}
即
\nopagebreak[4]
\[(1+x)^p>1+px .\]

\Needspace{5\baselineskip}
\textbf{（2）} 令
\nopagebreak[4]
\[\alpha=c^{\frac1p}>0,\qquad t_n=\frac{\alpha}{a_n}.\]
\Needspace{5\baselineskip}
由 $a_1>\alpha$ 得 $0<t_1<1$。假设 $0<t_n<1$，即 $a_n>\alpha$。把递推式改写为
\nopagebreak[4]
\[a_{n+1}=a_n\left(\frac{p-1}p+\frac{c}{p}a_n^{-p}\right)
=a_n\cdot\frac{(p-1)+t_n^p}{p},\]
其中用了 $c=\alpha^p$ 与 $t_n^p=\dfrac{\alpha^p}{a_n^p}$。

\Needspace{5\baselineskip}
先证 $a_{n+1}<a_n$。由 $0<t_n<1$ 得 $t_n^p<1$，故
\nopagebreak[4]
\[\frac{(p-1)+t_n^p}{p}<\frac{(p-1)+1}{p}=1,\]
从而 $a_{n+1}<a_n$。

\Needspace{5\baselineskip}
再证 $a_{n+1}>\alpha$。由 $a_n=\dfrac{\alpha}{t_n}$，
\nopagebreak[4]
\[a_{n+1}>\alpha
\iff \frac{(p-1)+t_n^p}{p}\cdot\frac{\alpha}{t_n}>\alpha
\iff (p-1)+t_n^p>pt_n .\]
\Needspace{5\baselineskip}
而
\nopagebreak[4]
\[(p-1)+t_n^p>pt_n
\iff t_n^p>1+p(t_n-1).\]
\Needspace{5\baselineskip}
令 $x=t_n-1$，则 $-1<x<0$ 且 $x\neq0$。由（1）的结论
\nopagebreak[4]
\[(1+x)^p>1+px,\]
\Needspace{5\baselineskip}
即
\nopagebreak[4]
\[t_n^p>1+p(t_n-1),\]
故 $a_{n+1}>\alpha$。

\Needspace{5\baselineskip}
于是
\nopagebreak[4]
\[\alpha<a_{n+1}<a_n,\]
\Needspace{5\baselineskip}
且 $0<t_{n+1}=\dfrac{\alpha}{a_{n+1}}<1$。由数学归纳法，对一切 $n\in\mathbf N^*$ 都有
\nopagebreak[4]
\[a_n>a_{n+1}>c^{\frac1p}.\]

\gkanswer{答案（1）：见解析版}

\gkanswer{答案（2）：见解析版}

\end{answer}